\documentclass{amsart}
% For dealing with references we use the comment environment
\usepackage{verbatim}
\newenvironment{reference}{\comment}{\endcomment}
%\newenvironment{reference}{}{}
\newenvironment{slogan}{\comment}{\endcomment}
\newenvironment{history}{\comment}{\endcomment}

% For commutative diagrams we use Xy-pic
\usepackage[all]{xy}

% We use 2cell for 2-commutative diagrams.
\xyoption{2cell}
\UseAllTwocells

% We use multicol for the list of chapters between chapters
\usepackage{multicol}

% This is generally recommended for better output
\usepackage{lmodern}
\usepackage[T1]{fontenc}

% For cross-file-references

% Package for hypertext links:
\usepackage{hyperref}

% For any local file, say "hello.tex" you want to link to please

% Theorem environments.
%
\theoremstyle{plain}
\newtheorem{theorem}[subsection]{Theorem}
\newtheorem{proposition}[subsection]{Proposition}
\newtheorem{lemma}[subsection]{Lemma}

\theoremstyle{definition}
\newtheorem{definition}[subsection]{Definition}
\newtheorem{example}[subsection]{Example}
\newtheorem{exercise}[subsection]{Exercise}
\newtheorem{situation}[subsection]{Situation}

\theoremstyle{remark}
\newtheorem{remark}[subsection]{Remark}
\newtheorem{remarks}[subsection]{Remarks}

\numberwithin{equation}{subsection}

% Macros
%
\def\lim{\mathop{\mathrm{lim}}\nolimits}
\def\colim{\mathop{\mathrm{colim}}\nolimits}
\def\Spec{\mathop{\mathrm{Spec}}}
\def\Hom{\mathop{\mathrm{Hom}}\nolimits}
\def\Ext{\mathop{\mathrm{Ext}}\nolimits}
\def\SheafHom{\mathop{\mathcal{H}\!\mathit{om}}\nolimits}
\def\SheafExt{\mathop{\mathcal{E}\!\mathit{xt}}\nolimits}
\def\Sch{\mathit{Sch}}
\def\Mor{\mathop{\mathrm{Mor}}\nolimits}
\def\Ob{\mathop{\mathrm{Ob}}\nolimits}
\def\Sh{\mathop{\mathit{Sh}}\nolimits}
\def\NL{\mathop{N\!L}\nolimits}
\def\CH{\mathop{\mathrm{CH}}\nolimits}
\def\proetale{{pro\text{-}\acute{e}tale}}
\def\etale{{\acute{e}tale}}
\def\QCoh{\mathit{QCoh}}
\def\Ker{\mathop{\mathrm{Ker}}}
\def\Im{\mathop{\mathrm{Im}}}
\def\Coker{\mathop{\mathrm{Coker}}}
\def\Coim{\mathop{\mathrm{Coim}}}

% Boxtimes
%
\DeclareMathSymbol{\boxtimes}{\mathbin}{AMSa}{"02}

%
% Macros for moduli stacks/spaces
%
\def\QCohstack{\mathcal{QC}\!\mathit{oh}}
\def\Cohstack{\mathcal{C}\!\mathit{oh}}
\def\Spacesstack{\mathcal{S}\!\mathit{paces}}
\def\Quotfunctor{\mathrm{Quot}}
\def\Hilbfunctor{\mathrm{Hilb}}
\def\Curvesstack{\mathcal{C}\!\mathit{urves}}
\def\Polarizedstack{\mathcal{P}\!\mathit{olarized}}
\def\Complexesstack{\mathcal{C}\!\mathit{omplexes}}
% \Pic is the operator that assigns to X its picard group, usage \Pic(X)
% \Picardstack_{X/B} denotes the Picard stack of X over B
% \Picardfunctor_{X/B} denotes the Picard functor of X over B
\def\Pic{\mathop{\mathrm{Pic}}\nolimits}
\def\Picardstack{\mathcal{P}\!\mathit{ic}}
\def\Picardfunctor{\mathrm{Pic}}
\def\Deformationcategory{\mathcal{D}\!\mathit{ef}}

\setlength{\emergencystretch}{3em}
\begin{document}
\title[Stacks Project, Tag 0A02]{308 --- Construction and universal property of henselization of pairs}
\maketitle
\noindent Copyright (C) 2005--2025 Johan de Jong. Stacks Project authors. Source: \href{https://stacks.math.columbia.edu/tag/0A02}{Tag 0A02}.

\noindent Editorial difficulty note (not author text). Difficulty H2: The proof constructs henselization as a directed colimit of etale neighbourhoods and explicitly verifies directedness, the residue quotient, and the henselian property. Its adjunction requires coherent uniqueness of sections proved by diagonal idempotents; existence is not delegated to a pre-existing henselization theorem..

\noindent Editorial provenance note (not author text). History: excerpt from revision \texttt{a0444\allowbreak 6e57e\allowbreak c1fbc\allowbreak 252a8\allowbreak 71afc\allowbreak ec775\allowbreak 2fb28\allowbreak 07b14}, more-algebra.tex, lines 2979--3084. Reference presentation was adapted; original-author context and core support blocks were retained or added. The mathematical wording is unchanged. Licensed under GNU FDL 1.2 or later; no invariant sections or cover texts. See ../licenses/COPYING. Context blocks reproduce author definitions and supporting results; they are not additional counted problems.

\begin{definition}
\label{categories-definition-directed}
We say that a diagram $M : \mathcal{I} \to \mathcal{C}$ is {\it directed},
or {\it filtered} if the following conditions hold:
\begin{enumerate}
\item the category $\mathcal{I}$ has at least one object,
\item for every pair of objects $x, y$ of $\mathcal{I}$
there exist an object $z$ and morphisms $x \to z$,
$y \to z$, and
\item for every pair of objects $x, y$ of $\mathcal{I}$
and every pair of morphisms $a, b : x \to y$ of $\mathcal{I}$
there exists a morphism $c : y \to z$ of $\mathcal{I}$
such that $M(c \circ a) = M(c \circ b)$ as morphisms in $\mathcal{C}$.
\end{enumerate}
We say that an index category $\mathcal{I}$ is {\it directed}, or
{\it filtered} if $\text{id} : \mathcal{I} \to \mathcal{I}$ is filtered
(in other words you erase the $M$ in part (3) above).
\end{definition}

\begin{lemma}
\label{lemma-henselization}
The inclusion functor
$$
\text{category of henselian pairs}
\longrightarrow
\text{category of pairs}
$$
has a left adjoint $(A, I) \mapsto (A^h, I^h)$.
\end{lemma}

\begin{proof}
Let $(A, I)$ be a pair. Consider the category $\mathcal{C}$ consisting
of \'etale ring maps $A \to B$ such that $A/I \to B/IB$ is an isomorphism.
We will show that the category $\mathcal{C}$ is directed and that
$A^h = \colim_{B \in \mathcal{C}} B$ with ideal $I^h = IA^h$ gives
the desired adjoint.

\medskip\noindent
We first prove that $\mathcal{C}$ is directed
(Categories, Definition \ref{categories-definition-directed}).
It is nonempty because $\text{id} : A \to A$ is an object.
If $B$ and $B'$ are two objects of $\mathcal{C}$, then
$B'' = B \otimes_A B'$ is an object of $\mathcal{C}$
(use Algebra, Lemma \href{https://stacks.math.columbia.edu/tag/00U2}{lemma etale (Tag 00U2)})
and there are morphisms $B \to B''$ and $B' \to B''$.
Suppose that $f, g : B \to B'$ are two maps between
objects of $\mathcal{C}$. Then a coequalizer is
$$
(B' \otimes_{f, B, g} B') \otimes_{(B' \otimes_A B')} B'
$$
which is \'etale over $A$ by
Algebra, Lemmas \href{https://stacks.math.columbia.edu/tag/00U2}{lemma etale (Tag 00U2)} and
\href{https://stacks.math.columbia.edu/tag/00U7}{lemma map between etale (Tag 00U7)}.
Thus the category $\mathcal{C}$ is directed.

\medskip\noindent
Since $B/IB = A/I$ for all objects $B$ of $\mathcal{C}$ we
see that $A^h/I^h = A^h/IA^h = \colim B/IB = \colim A/I = A/I$.

\medskip\noindent
Next, we show that $A^h = \colim_{B \in \mathcal{C}} B$ with
$I^h = IA^h$ is a henselian pair. To do this we will verify
condition (2) of Lemma \href{https://stacks.math.columbia.edu/tag/09XI}{lemma characterize henselian pair (Tag 09XI)}.
Namely, suppose given an \'etale ring map $A^h \to A'$
and $A^h$-algebra map $\sigma : A' \to A^h/I^h$. Then there exists a
$B \in \mathcal{C}$ and an \'etale ring map $B \to B'$ such that
$A' = B' \otimes_B A^h$. See Algebra, Lemma \href{https://stacks.math.columbia.edu/tag/00U2}{lemma etale (Tag 00U2)}.
Since $A^h/I^h = A/I$, the map $\sigma$ induces an $A$-algebra
map $s : B' \to A/I$. Then $B'/IB' = A/I \times C$ as $A/I$-algebra,
where $C$ is the kernel of the map $B'/IB' \to A/I$ induced by $s$.
Let $g \in B'$ map to $(1, 0) \in A/I \times C$. Then $B \to B'_g$
is \'etale and $A/I \to B'_g/IB'_g$ is an isomorphism, i.e.,
$B'_g$ is an object of $\mathcal{C}$. Thus we obtain a canonical
map $B'_g \to A^h$ such that
$$
\vcenter{
\xymatrix{
B'_g \ar[r] & A^h \\
B \ar[u] \ar[ur]
}
}
\quad\text{and}\quad
\vcenter{
\xymatrix{
B' \ar[r] \ar[rrd]_s & B'_g \ar[r] & A^h \ar[d] \\
& & A/I
}
}
$$
commute. This induces a map $A' = B' \otimes_B A^h \to A^h$
compatible with $\sigma$ as desired.

\medskip\noindent
Let $(A, I) \to (A', I')$ be a morphism of pairs with $(A', I')$ henselian.
We will show there is a unique factorization $A \to A^h \to A'$ which will
finish the proof. Namely, for each $A \to B$ in $\mathcal{C}$
the ring map $A' \to B' = A' \otimes_A B$ is \'etale and induces
an isomorphism $A'/I' \to B'/I'B'$. Hence there is a section
$\sigma_B : B' \to A'$ by Lemma \href{https://stacks.math.columbia.edu/tag/09XI}{lemma characterize henselian pair (Tag 09XI)}.
Given a morphism $B_1 \to B_2$ in $\mathcal{C}$ we claim the diagram
$$
\xymatrix{
B'_1 \ar[rr] \ar[rd]_{\sigma_{B_1}} & &
B'_2 \ar[ld]^{\sigma_{B_2}} \\
& A'
}
$$
commutes. This follows once we prove that for every $B$ in $\mathcal{C}$
the section $\sigma_B$ is the unique $A'$-algebra map $B' \to A'$.
We have $B' \otimes_{A'} B' = B' \times R$ for some ring $R$, see
Algebra, Lemma \href{https://stacks.math.columbia.edu/tag/02FL}{lemma diagonal unramified (Tag 02FL)}. In our case
$R/I'R = 0$ as $B'/I'B' = A'/I'$. Thus given two $A'$-algebra maps
$\sigma_B, \sigma_B' : B' \to A'$ then
$e = (\sigma_B \otimes \sigma_B')(0, 1) \in A'$
is an idempotent contained in $I'$. We conclude that $e = 0$
by Lemma \href{https://stacks.math.columbia.edu/tag/09XF}{lemma idempotents determined modulo radical (Tag 09XF)}.
Hence $\sigma_B = \sigma_B'$ as desired.
Using the commutativity we obtain
$$
A^h = \colim_{B \in \mathcal{C}} B \to
\colim_{B \in \mathcal{C}} A' \otimes_A B \xrightarrow{\colim \sigma_B} A'
$$
as desired. The uniqueness of the maps $\sigma_B$ also guarantees that
this map is unique. Hence $(A, I) \mapsto (A^h, I^h)$ is the desired adjoint.
\end{proof}
\end{document}
